% Modul Belajar 3, dihasilkan oleh tools/build.py dari tools/konten/p03.py
\documentclass[a4paper]{article}
\usepackage{modulITERA}
\begin{document}
\Sampul{3}{Operator dan Precedence}{Sisa bagi, urutan operasi, penugasan majemuk, serta operator relasional dan logika}{CPMK0607 (menerapkan) dan CPMK0608 (menjelaskan) pada konsep dasar algoritma dan sistem komputer}{sekitar 3 jam belajar mandiri, ditambah 150 menit di kelas}{\texttt{02 Hands-on/P3 - Operator - Hands-on/}}
\Bagian{Modul 3}{Peta belajar}
Ikuti urutan ini dari atas ke bawah. Setiap bagian memakai apa yang sudah dibahas di bagian sebelumnya.
\par\vspace{4pt}\noindent{\fontsize{9.5pt}{13pt}\selectfont
\begin{tabular}{@{}>{\raggedright\arraybackslash}p{0.140\textwidth}>{\raggedright\arraybackslash}p{0.840\textwidth}@{}}
\kepala{Urutan} & \kepala{Bagian}\\
1 & Tentang modul ini\\\garis
2 & Masalah pembuka: Durasi lagu di playlist\\\garis
3 & Langkah 1: Sisa bagi dengan \%\\\garis
4 & Langkah 2: Urutan pengerjaan operator\\\garis
5 & Langkah 3: Menelusuri urutan operasi\\\garis
6 & Langkah 4: Penugasan majemuk dan increment\\\garis
7 & Langkah 5: Membandingkan dan menggabungkan kondisi\\\garis
8 & Langkah 6: Konversi tipe dan fungsi matematika\\\garis
9 & Pesan error yang sering muncul\\\garis
10 & Latihan mandiri\\\garis
11 & Rangkuman\\\garis
12 & Cek pemahaman\\\garis
13 & Exit Ticket dan tugas\\\garis
14 & Bacaan lanjut dan persiapan minggu depan\\
\garisdasar
\end{tabular}}\par\vspace{6pt}
\Bagian{Pembuka}{Tentang modul ini}
Modul ini mendampingi Pertemuan 3. Kalian sudah bisa menyimpan dan membaca data. Minggu ini kalian belajar mengolahnya dengan operator yang lebih lengkap, dan yang paling penting, memahami urutan C++ mengerjakan sebuah ekspresi.

Isinya sama dengan slide di kelas, tetapi ditulis lebih pelan dan lebih lengkap, supaya kalian bisa mengulang sendiri di rumah tanpa harus menebak maksud slide.

\Sub{Setelah menyelesaikan modul ini, kalian bisa}
\par\vspace{4pt}\noindent{\fontsize{9.5pt}{13pt}\selectfont
\begin{tabular}{@{}>{\raggedright\arraybackslash}p{0.070\textwidth}>{\raggedright\arraybackslash}p{0.680\textwidth}>{\raggedright\arraybackslash}p{0.230\textwidth}@{}}
\kepala{No} & \kepala{Kemampuan} & \kepala{Dibahas di}\\
1 & Menulis program yang memakai operator aritmetika termasuk \texttt{\%}, operator penugasan majemuk, increment, serta operator relasional dan logika untuk menyelesaikan perhitungan (Sub-CPMK 3.1, CPMK0607) & Langkah 1 sampai 6\\\garis
2 & Menjelaskan urutan evaluasi sebuah ekspresi berdasarkan precedence dan asosiativitas, serta nilai yang dihasilkan pada setiap langkah (Sub-CPMK 3.2, CPMK0608) & kotak Tebak dulu, Latihan 3\\
\garisdasar
\end{tabular}}\par\vspace{6pt}
Karena setiap instrumen penilaian dibagi rata antara Bagian A (menulis kode) dan Bagian B (menelusuri dan menjelaskan), yang dinilai bukan hanya program kalian jalan, tetapi juga kalian bisa menjelaskan mengapa program itu jalan.

\Sub{Cara memakai modul ini}
Setiap langkah punya pola yang sama: penjelasan konsep, kadang contoh dari kehidupan sehari-hari, lalu kode yang bisa langsung kalian ketik. Sesudah kode ada kotak \textbf{Tebak dulu}. Tulis tebakan kalian di kertas sebelum membaca \textbf{Jawaban} di bawahnya. Kebiasaan menebak lalu membuktikan inilah yang membuat konsep menempel, bukan sekadar membaca.

\begin{Penting}[Ketik, jangan salin]
Ketik ulang setiap contoh di GDB Online atau editor kalian, jangan disalin-tempel. Saat mengetik, kalian akan membuat salah ketik, membaca pesan error, dan memperbaikinya. Itu bagian dari belajar.
\end{Penting}
\Sub{Yang perlu disiapkan}
Peramban dengan akses ke onlinegdb.com, atau g++ di komputer kalian sendiri. Kertas dan alat tulis untuk menebak keluaran dan membuat trace table. Folder hands-on pertemuan ini dari repo mata kuliah.

\Bagian{Pembuka}{Masalah pembuka: Durasi lagu di playlist}
Aplikasi musik menyimpan durasi playlist dalam detik, misalnya 3725 detik. Pengguna ingin melihatnya sebagai jam, menit, dan detik.

Pembagian biasa memberi 1.03 jam, yang tidak membantu. Kita perlu cara mengambil bagian bulat dan sisanya.

\begin{MKeluaran}[title={Tampilan yang diinginkan}]
Durasi: 3725 detik
= 1 jam 2 menit 5 detik
\end{MKeluaran}
\begin{Tebak}[Coba pikirkan]
Bagaimana kalian menghitung ini dengan tangan? Operasi apa yang dipakai untuk mendapat angka 5 di akhir?
\end{Tebak}
\begin{Jawaban}[Cara memecahnya]
Pecah: Ambil jam, lalu sisa detiknya, lalu menit, lalu sisa lagi. Rumuskan: Jam adalah hasil bagi bulat; sisa adalah sisa bagi. Terjemahkan: Pakai \texttt{/} untuk hasil bagi bulat dan \texttt{\%} untuk sisa.
\end{Jawaban}
\Bagian{Langkah 1}{Sisa bagi dengan \%}
Operator \texttt{\%} memberi sisa pembagian bulat: \texttt{17 \% 5} bernilai 2 karena 17 = 3 × 5 + 2. Pasangan \texttt{/} dan \texttt{\%} sangat berguna untuk memecah satuan: detik ke jam, rupiah ke lembar uang, angka ke digit-digitnya.

\texttt{\%} hanya berlaku untuk bilangan bulat. Bila salah satu operand \texttt{double}, compiler menolak.

\begin{MKode}{modulo.cpp}
int total = 3725;
int jam = total / 3600;
int menit = total % 3600 / 60;
int detik = total % 60;
cout << jam << " " << menit << " " << detik << endl;
cout << 17 % 5 << " " << 20 % 4 << endl;
\end{MKode}
\begin{Tebak}[]
Sebelum menjalankan program, tulis keluarannya di kertas, karakter demi karakter.
\end{Tebak}
\begin{Jawaban}[]
\begin{MKeluaran}[]
1 2 5
2 0
\end{MKeluaran}
\texttt{\%} hanya untuk bilangan bulat. \texttt{a \% b == 0} berarti a habis dibagi b.
\end{Jawaban}
\Bagian{Langkah 2}{Urutan pengerjaan operator}
Bila sebuah ekspresi berisi beberapa operator, C++ mengerjakannya menurut tingkat (precedence). Operator setingkat dikerjakan menurut arahnya (asosiativitas). Tabel ini cukup untuk seluruh semester.

\par\vspace{4pt}\noindent{\fontsize{9.5pt}{13pt}\selectfont
\begin{tabular}{@{}>{\raggedright\arraybackslash}p{0.120\textwidth}>{\raggedright\arraybackslash}p{0.316\textwidth}>{\raggedright\arraybackslash}p{0.196\textwidth}>{\raggedright\arraybackslash}p{0.348\textwidth}@{}}
\kepala{Tingkat} & \kepala{Operator} & \kepala{Arah} & \kepala{Contoh}\\
1 (pertama) & \texttt{()} & dalam ke luar & \texttt{(2 + 3) * 4} = 20\\\garis
2 & \texttt{++ -- !} dan cast & kanan ke kiri & \texttt{!true} = false\\\garis
3 & \texttt{* / \%} & kiri ke kanan & \texttt{20 / 4 * 2} = 10\\\garis
4 & \texttt{+ -} & kiri ke kanan & \texttt{10 - 4 - 3} = 3\\\garis
5 & \texttt{< <= > >=} lalu \texttt{== !=} & kiri ke kanan & \texttt{3 < 5 == true}\\\garis
6 & \texttt{\&\&} lalu \texttt{||} & kiri ke kanan & \texttt{a || b \&\& c}\\\garis
7 (terakhir) & \texttt{= += -= *= /= \%=} & kanan ke kiri & \texttt{a = b = 0}\\
\garisdasar
\end{tabular}}\par\vspace{6pt}
\begin{Penting}[]
Ragu? Tambahkan kurung. Kurung tidak membuat program lambat, tetapi membuatnya mudah dibaca.
\end{Penting}
\Bagian{Langkah 3}{Menelusuri urutan operasi}
Cara menelusuri ekspresi: cari operator dengan tingkat tertinggi, kerjakan, tulis ulang ekspresinya, ulangi. Untuk \texttt{7 + 10 \% 4 * 2}, \texttt{\%} dan \texttt{*} setingkat dan dikerjakan dari kiri: \texttt{10 \% 4} = 2, lalu \texttt{2 * 2} = 4, terakhir \texttt{7 + 4} = 11.

\begin{MKode}{urutan.cpp}
cout << 2 + 3 * 4 << endl;
cout << (2 + 3) * 4 << endl;
cout << 20 / 4 * 2 << endl;
cout << 10 - 4 - 3 << endl;
cout << 7 + 10 % 4 * 2 << endl;
\end{MKode}
\begin{Tebak}[]
Sebelum menjalankan program, tulis keluarannya di kertas, karakter demi karakter.
\end{Tebak}
\begin{Jawaban}[]
\begin{MKeluaran}[]
14
20
10
3
11
\end{MKeluaran}
Baris terakhir: \texttt{10 \% 4} = 2, lalu \texttt{2 * 2} = 4, lalu \texttt{7 + 4} = 11.
\end{Jawaban}
\begin{Penting}[]
Di Bagian B, tuliskan langkahnya satu per satu, bukan hanya hasil akhirnya.
\end{Penting}
\Bagian{Langkah 4}{Penugasan majemuk dan increment}
\texttt{x += 3} adalah singkatan \texttt{x = x + 3}, begitu juga \texttt{-=}, \texttt{*=}, \texttt{/=}, dan \texttt{\%=}. \texttt{x++} dan \texttt{x--} menambah atau mengurangi satu.

Perbedaan \texttt{x++} dan \texttt{++x} hanya terlihat bila dipakai di dalam ekspresi lain. \texttt{int a = x++} menyalin nilai lama ke \texttt{a}, baru \texttt{x} bertambah. \texttt{int b = ++x} menambah \texttt{x} dulu, lalu menyalin nilai barunya.

\begin{MKode}{majemuk.cpp}
int x = 5;
x += 3;
x *= 2;
x--;
int a = x++;
int b = ++x;
cout << a << " " << b << " " << x << endl;
\end{MKode}
\begin{Tebak}[]
Sebelum menjalankan program, tulis keluarannya di kertas, karakter demi karakter.
\end{Tebak}
\begin{Jawaban}[]
\begin{MKeluaran}[]
15 17 17
\end{MKeluaran}
\texttt{x++} memakai nilai lama dulu, lalu menambah. \texttt{++x} menambah dulu, lalu memakai nilai baru.
\end{Jawaban}
\Bagian{Langkah 5}{Membandingkan dan menggabungkan kondisi}
Operator relasional (\texttt{< <= > >= == !=}) menghasilkan nilai \texttt{bool}. Operator logika menggabungkannya: \texttt{\&\&} benar bila keduanya benar, \texttt{||} benar bila salah satu benar, \texttt{!} membalik nilai.

Hasil inilah yang minggu depan dipakai \texttt{if} untuk memilih jalan program. \texttt{boolalpha} membuat \texttt{cout} menampilkan \texttt{true} dan \texttt{false}, bukan 1 dan 0.

\begin{MKode}{logika.cpp}
int nilai = 78, hadir = 12;
bool lulus = nilai >= 60;
bool aktif = hadir >= 11;
cout << boolalpha;
cout << lulus << " " << aktif << endl;
cout << (lulus && aktif) << endl;
cout << (nilai >= 85 || hadir == 14) << endl;
cout << !lulus << endl;
\end{MKode}
\begin{Tebak}[]
Sebelum menjalankan program, tulis keluarannya di kertas, karakter demi karakter.
\end{Tebak}
\begin{Jawaban}[]
\begin{MKeluaran}[]
true true
true
false
false
\end{MKeluaran}
\texttt{==} membandingkan, \texttt{=} menugaskan. Tanpa \texttt{boolalpha}, true dan false tampil sebagai 1 dan 0.
\end{Jawaban}
\Bagian{Langkah 6}{Konversi tipe dan fungsi matematika}
Cast mengubah tipe sebuah nilai: \texttt{(int) 3.99} menjadi 3 karena bagian pecahan dipotong, bukan dibulatkan. Untuk membulatkan, pakai \texttt{round}.

Pustaka \texttt{<cmath>} menyediakan fungsi matematika seperti \texttt{sqrt} (akar), \texttt{pow} (pangkat), \texttt{round}, \texttt{floor}, dan \texttt{ceil}.

\begin{MKode}{konversi.cpp}
double d = 3.99;
int i = (int) d;
cout << i << endl;
cout << (double) 7 / 2 << endl;
cout << sqrt(16.0) << " " << pow(2, 10) << endl;
cout << round(2.5) << " " << floor(2.7) << endl;
\end{MKode}
\begin{Tebak}[]
Sebelum menjalankan program, tulis keluarannya di kertas, karakter demi karakter.
\end{Tebak}
\begin{Jawaban}[]
\begin{MKeluaran}[]
3
3.5
4 1024
3 2
\end{MKeluaran}
Cast ke \texttt{int} memotong, tidak membulatkan. Fungsi matematika perlu \texttt{\#include <cmath>}.
\end{Jawaban}
\Bagian{Waspada}{Pesan error yang sering muncul}
Semua pesan di tabel ini dihasilkan g++ dengan opsi \texttt{-Wall}, sama seperti yang dipakai \texttt{cek.sh}.
\par\vspace{4pt}\noindent{\fontsize{9.5pt}{13pt}\selectfont
\begin{tabular}{@{}>{\raggedright\arraybackslash}p{0.240\textwidth}>{\raggedright\arraybackslash}p{0.270\textwidth}>{\raggedright\arraybackslash}p{0.470\textwidth}@{}}
\kepala{Kesalahan} & \kepala{Contoh} & \kepala{Pesan pertama dari g++}\\
\% pada bilangan pecahan & \texttt{5.5 \% 2} & \texttt{error: invalid operands of types 'double' and 'int' to binary 'operator\%'}\\\garis
Ruas kiri bukan variabel & \texttt{5 = x;} & \texttt{error: lvalue required as left operand of assignment}\\\garis
Pembagian dengan nol & \texttt{int y = 10 / 0;} & \texttt{warning: division by zero}\\\garis
Perbandingan tanpa kurung di cout & \texttt{cout << 3 > 2;} & \texttt{error: no match for 'operator>'}\\\garis
== yang tidak dipakai & \texttt{x == 5;} & \texttt{warning: statement has no effect}\\
\garisdasar
\end{tabular}}\par\vspace{6pt}
Kesalahan precedence biasanya tidak memunculkan pesan apa pun. Program jalan, hasilnya salah. Itu sebabnya latihan tantangan minggu ini adalah menelusuri dan memperbaiki rumus.

\Bagian{Latihan}{Latihan mandiri}
Latihan ini sama dengan hands-on di kelas. Semua file ada di \texttt{02 Hands-on/P3 - Operator - Hands-on/latihan/}. Kerjakan berurutan, lalu cocokkan dengan \texttt{expected/} atau jalankan \texttt{bash cek.sh}.
\par\vspace{4pt}\noindent{\fontsize{9.5pt}{13pt}\selectfont
\begin{tabular}{@{}>{\raggedright\arraybackslash}p{0.070\textwidth}>{\raggedright\arraybackslash}p{0.360\textwidth}>{\raggedright\arraybackslash}p{0.150\textwidth}>{\raggedright\arraybackslash}p{0.400\textwidth}@{}}
\kepala{No} & \kepala{File} & \kepala{Tingkat} & \kepala{Yang dilatih}\\
1 & \texttt{handson1-waktu.cpp} & Dasar & \texttt{/} dan \texttt{\%} untuk memecah satuan\\\garis
2 & \texttt{handson2-kembalian.cpp} & Menengah & \texttt{\%=} dan pembagian bulat berulang\\\garis
3 & \texttt{handson3-rumus.cpp} & Tantangan & telusuri precedence, perbaiki tiga rumus\\\garis
+ & \texttt{bonus-genap.cpp} & Pengayaan & relasional, \texttt{\&\&}, \texttt{boolalpha}\\
\garisdasar
\end{tabular}}\par\vspace{6pt}
\Sub{Latihan 1: handson1-waktu.cpp}
Ubah jumlah detik menjadi jam, menit, dan detik. Masukan uji: \texttt{3725}.

\Sub{Latihan 2: handson2-kembalian.cpp}
Pecah uang kembalian menjadi lembar 50000 sampai 1000. Masukan uji: \texttt{100000 63000}.

\Sub{Latihan 3: handson3-rumus.cpp}
Tiga rumus salah karena urutan operasi dan pembagian bulat. Telusuri, perbaiki, dan jelaskan di blok \texttt{PENJELASAN}.

\Sub{Bonus: bonus-genap.cpp}
Tampilkan sifat sebuah bilangan sebagai nilai \texttt{bool} tanpa \texttt{if}.

\Sub{Latihan menelusuri}
\begin{MKode}{tebak.cpp}
int a = 9, b = 4;
int c = a % b + a / b * 2;
a += c;
b = a-- - b;
cout << c << " " << a << " " << b << endl;
cout << (a > b && c != 5) << endl;
\end{MKode}
\begin{Tebak}[]
Tulis keluarannya di kertas tanpa menjalankan program. Buat trace table bila perlu.
\end{Tebak}
\begin{Jawaban}[]
\begin{MKeluaran}[]
5 13 10
0
\end{MKeluaran}
\texttt{c} = 1 + 2 × 2 = 5. \texttt{a} menjadi 14. \texttt{b} dihitung dengan nilai lama \texttt{a} (14 − 4 = 10), lalu \texttt{a} turun ke 13. Karena \texttt{c} = 5, kondisi terakhir bernilai 0.
\end{Jawaban}
\Bagian{Penutup}{Rangkuman}
\par\vspace{4pt}\noindent{\fontsize{9.5pt}{13pt}\selectfont
\begin{tabular}{@{}>{\raggedright\arraybackslash}p{0.320\textwidth}>{\raggedright\arraybackslash}p{0.660\textwidth}@{}}
\kepala{Konsep} & \kepala{Yang perlu diingat}\\
Sisa bagi dengan \% & \texttt{\%} hanya untuk bilangan bulat.\\\garis
Urutan pengerjaan operator & Ragu?\\\garis
Menelusuri urutan operasi & Baris terakhir: \texttt{10 \% 4} = 2, lalu \texttt{2 * 2} = 4, lalu \texttt{7 + 4} = 11.\\\garis
Penugasan majemuk dan increment & \texttt{x++} memakai nilai lama dulu, lalu menambah.\\\garis
Membandingkan dan menggabungkan kondisi & \texttt{==} membandingkan, \texttt{=} menugaskan.\\\garis
Konversi tipe dan fungsi matematika & Cast ke \texttt{int} memotong, tidak membulatkan.\\
\garisdasar
\end{tabular}}\par\vspace{6pt}
\Bagian{Penutup}{Cek pemahaman}
Jawab tanpa menjalankan program, lalu periksa dengan menjalankannya.
\begin{enumerate}
\item Berapa hasil \texttt{25 \% 7} dan \texttt{25 / 7}?
\item Berapa nilai \texttt{3 + 4 * 2 - 6 / 3}?
\item Sesudah \texttt{int x = 4; int y = x++ * 2;}, berapa \texttt{x} dan \texttt{y}?
\item Apa nilai \texttt{!(5 > 3) || 2 == 2}?
\item Mengapa \texttt{(int) 2.9} bernilai 2?
\end{enumerate}
\begin{Jawaban}[]
\begin{enumerate}[leftmargin=16pt]
\item 4 dan 3.
\item 9 (4 × 2 = 8, 6 / 3 = 2, 3 + 8 − 2).
\item \texttt{x} = 5, \texttt{y} = 8.
\item true, karena ruas kanan benar.
\item Cast ke \texttt{int} memotong bagian pecahan, tidak membulatkan.
\end{enumerate}
\end{Jawaban}
\Bagian{Penilaian}{Exit Ticket 3}
Dikerjakan di 25 menit terakhir kelas, sendiri, tanpa AI, internet, atau diskusi. Exit Ticket 3 adalah satu dari 14 Exit Ticket; rata-ratanya menjadi komponen berbobot 25\%.
\par\vspace{4pt}\noindent{\fontsize{9.5pt}{13pt}\selectfont
\begin{tabular}{@{}>{\raggedright\arraybackslash}p{0.080\textwidth}>{\raggedright\arraybackslash}p{0.600\textwidth}>{\raggedright\arraybackslash}p{0.100\textwidth}>{\raggedright\arraybackslash}p{0.200\textwidth}@{}}
\kepala{Soal} & \kepala{Isi} & \kepala{Poin} & \kepala{CPMK}\\
A1 & Tulis program yang memecah sebuah bilangan tiga digit menjadi ratusan, puluhan, dan satuan & 25 & CPMK0607\\\garis
A2 & Tulis ekspresi logika untuk syarat yang diberikan dan tampilkan hasilnya & 25 & CPMK0607\\\garis
B1 & Tunjukkan langkah evaluasi sebuah ekspresi campuran sesuai precedence & 20 & CPMK0608\\\garis
B2 & Telusuri nilai variabel setelah penugasan majemuk dan increment & 20 & CPMK0608\\\garis
B3 & Jelaskan beda \texttt{x++} dan \texttt{++x} dengan contoh & 10 & CPMK0608\\
\garisdasar
\end{tabular}}\par\vspace{6pt}
\Bagian{Penilaian}{Tugas 3: Kalkulator Lengkap dengan Validasi}
Kumpulkan sebelum Pertemuan 4 lewat kanal pengumpulan yang diumumkan dosen.

\Sub{Bagian A: program (50 poin)}
Buat program kalkulator yang membaca dua angka desimal, lalu menampilkan hasil berikut.
\begin{enumerate}
\item Masukan dua angka bertipe \texttt{double}
\item Hasil \texttt{+}, \texttt{-}, \texttt{*}, dan \texttt{/}
\item Sisa bagi \texttt{\%} dari bagian bulat kedua angka (cast ke \texttt{int})
\item Sebelum membagi, periksa apakah pembagi nol; tampilkan pesan dengan operator ternary \texttt{?:}
\item Nilai \texttt{bool}: apakah angka pertama lebih besar
\item Nilai \texttt{bool}: apakah hasil penjumlahan > 100 dan hasil kali > 50
\end{enumerate}
\begin{Contoh}[Bonus, tidak wajib]
Tampilkan apakah bagian bulat masing-masing angka ganjil atau genap.
\end{Contoh}
\Sub{Bagian B: penjelasan (50 poin)}
Komentar maksimal 150 kata: urutan evaluasi ekspresi syarat terakhir langkah demi langkah, alasan perlu cast sebelum \texttt{\%}, dan satu kesalahan precedence atau pembagian yang kalian temui.

\Sub{Rubrik Tugas 3}
\par\vspace{4pt}\noindent{\fontsize{8.5pt}{11.5pt}\selectfont
\begin{tabular}{@{}>{\raggedright\arraybackslash}p{0.190\textwidth}>{\raggedright\arraybackslash}p{0.070\textwidth}>{\raggedright\arraybackslash}p{0.200\textwidth}>{\raggedright\arraybackslash}p{0.180\textwidth}>{\raggedright\arraybackslash}p{0.170\textwidth}>{\raggedright\arraybackslash}p{0.170\textwidth}@{}}
\kepala{Kriteria} & \kepala{Poin} & \kepala{Sangat baik 85–100} & \kepala{Baik 70–84} & \kepala{Cukup 55–69} & \kepala{Kurang <55}\\
A1 Program berjalan & 20 & tanpa error dan peringatan & ada peringatan & perlu satu perbaikan & tidak terkompilasi\\\garis
A2 Kelengkapan operasi & 15 & semua operasi dan validasi & satu operasi kurang & dua kurang & tiga atau lebih kurang\\\garis
A3 Validasi dan logika & 15 & pembagi nol dan dua bool tepat & satu salah & dua salah & tidak ada validasi\\\garis
B1 Langkah evaluasi & 30 & urutan tepat dan beralasan & satu langkah keliru & hanya hasil akhir & tidak ada atau disalin\\\garis
B2 Cerita error dan pengujian & 20 & pesan, sebab, perbaikan, uji & pesan tidak dikutip & hanya perbaikan & tidak ada\\
\garisdasar
\end{tabular}}\par\vspace{6pt}
Nilai kriteria = poin × skor level. Contoh: A1 di level Baik dengan skor 80 memberi 20 × 0,80 = 16 poin.

\Bagian{Penutup}{Bacaan lanjut dan persiapan minggu depan}
\begin{enumerate}
\item Deitel, P. dan Deitel, H. C++ How to Program, edisi ke-10, Bab 2 dan 4.
\item learncpp.com, Bab 6 (operator).
\item cppreference.com, C++ Operator Precedence.
\item Slide \texttt{01 Slide/P3 Operator dan Precedence.pdf} dan hands-on di \texttt{02 Hands-on/P3 - Operator - Hands-on/}.
\end{enumerate}
\begin{Penting}[Minggu depan]
Pertemuan 4 membahas percabangan dengan \texttt{if}, \texttt{else if}, dan \texttt{switch}, memakai ekspresi relasional dan logika dari minggu ini.
\end{Penting}
\end{document}
